%; whizzy chapter
% -initex iniptex -latex platex -format platex -bibtex jbibtex -fmt fmt
% 以上 whizzytex を使用する場合の設定。

%     Tokyo Debian Meeting resources
%     Copyright (C) 2006 Junichi Uekawa

%     This program is free software; you can redistribute it and/or modify
%     it under the terms of the GNU General Public License as published by
%     the Free Software Foundation; either version 2 of the License, or
%     (at your option) any later version.

%     This program is distributed in the hope that it will be useful,
%     but WITHOUT ANY WARRANTY; without even the implied warranty of
%     MERCHANTABILITY or FITNESS FOR A PARTICULAR PURPOSE.  See the
%     GNU General Public License for more details.

%     You should have received a copy of the GNU General Public License
%     along with this program; if not, write to the Free Software
%     Foundation, Inc., 51 Franklin St, Fifth Floor, Boston, MA  02110-1301 USA

%   Pdf作成手順
% dvipdfmx debianmeetingresume200601.dvi
%  preview (shell-command (concat "xpdf " (replace-regexp-in-string "tex$" "pdf"(buffer-file-name)) "&"))

%%ここからヘッダ開始。

\documentclass[mingoth,a4paper]{jsarticle}
\usepackage[dvipdfmx]{graphicx}
\usepackage{fancybox}
\usepackage{longtable}
\usepackage{ascmac}	% 囲み (screen,itembox)
\usepackage{fancyvrb}   % 囲み Verbatim のために必要
\usepackage[dvipdfmx]{hyperref}
\usepackage{url}

%http://www.naney.org/diki/dk/hyperref.html
%日本語EUC系環境の時
\AtBeginDvi{\special{pdf:tounicode EUC-UCS2}}
%シフトJIS系環境の時
%\AtBeginDvi{\special{pdf:tounicode 90ms-RKSJ-UCS2}}

%% spacing の設定をする。外枠を減らす。
\setlength\headheight{0mm}
\setlength\topmargin{-20mm}
\setlength\headsep{0mm}
\setlength\topskip{3mm}
\setlength\maxdepth{4pt}
\setlength\columnsep{6mm}
\setlength\textheight{252mm}
\setlength\topmargin{-5mm}
\setlength\textwidth{170mm}
\setlength\oddsidemargin{-5mm}
\setlength\evensidemargin{-5mm}

% commandline環境を定義。画面入出力についてはcommandline環境
% で表記する
\newenvironment{commandline}%
{\VerbatimEnvironment
  \begin{Sbox}\begin{minipage}{15cm}\begin{fontsize}{7.3}{7.3} \begin{BVerbatim}}%
{\end{BVerbatim}\end{fontsize}\end{minipage}\end{Sbox}
  \setlength{\fboxsep}{8pt}\fbox{\TheSbox}}


%%% start of santaku
\makeatletter
\newwrite\tf@jqz
\immediate\openout\tf@jqz\jobname.jqz\relax
\makeatother
\newcounter{santakucounter}
\newcommand{\santaku}[5]{%
\addtocounter{santakucounter}{1}

\addtocontents{jqz}{\arabic{santakucounter}. #5\\}
\nopagebreak 問題\arabic{santakucounter}. 
#1\\
\nopagebreak□ A #2\\
\nopagebreak□ B #3\\
\nopagebreak□ C #4
\pagebreak[1]
\hspace{1cm}
\\

}
%%% end of santaku

\newcommand{\emptyspace}{(\underline{\hspace{1cm}})}

\newcommand{\subsubsubsection}[1]{%
\vspace{1zw}{\bf #1}\\}

% sectionをセンタリングする
\makeatletter
  \renewcommand{\section}{\@startsection{section}{1}{\z@}%
    {\Cvs \@plus.5\Cdp \@minus.2\Cdp}% 前アキ
    {.5\Cvs \@plus.3\Cdp}% 後アキ
    {\normalfont\Large\headfont\raggedright\centering}} % style
\makeatother

% section の代わりの環境
\newcommand{\dancersection}[2]{%
\newpage
東京エリアDebian勉強会 2006
\hrule
\vspace{0.5mm}
\hrule
\hfill{}\includegraphics[width=3cm]{image200502/openlogo-nd.eps}\\
\vspace{-4cm}
\begin{center}
  \section{#1}
\end{center}
\hfill{}#2\hspace{3cm}\space\\
\hrule
\hrule
\vspace{1cm}
}

% for dancerj
\newcommand{\fgref}[1]{図\ref{#1}}
\newcommand{\tbref}[1]{表\ref{#1}}


\begin{document}

\begin{titlepage}

% 毎月変更する部分, 本文の末尾も修正することをわすれずに
\title{
 第12回 東京エリア Debian 勉強会\\事前資料}
\date{2006年1月21日}
\author{Debian勉強会会場係 上川 純一\thanks{Debian Project Official Developer}} 
\maketitle
\thispagestyle{empty}
\end{titlepage}

\newpage
\tableofcontents

\dancersection{Introduction To Debian 勉強会}{上川 純一}

今月のDebian勉強会へようこそ。
これからDebianのあやしい世界に入るという方も、すでにどっぷりとつかってい
るという方も、月に一回Debianについて語りませんか？

目的として下記の二つを考えています。

\begin{itemize}
 \item メールではよみとれない、もしくはよみとってられないような情報を情
       報共有する場をつくる
 \item まとまっていないDebianを利用する際の情報をまとめて、ある程度の塊と
       して出してみる
\end{itemize}

また、東京にはLinuxの勉強会はたくさんありますので、Debianに限定した勉強
会にします。Linuxの基本的な利用方法などが知りたい方は、他でがんばってくださ
い。
Debianの勉強会ということで究極的には参加者全員がDebian Packageを
がりがりと作りながらスーパーハッカーになれるような姿を妄想しています。

Debianをこれからどうするという能動的な展開への土台としての空間を提供し、
情報の共有をしたい、というのが目的です。
次回は違うこと言ってるかもしれませんが、御容赦を。

\subsection{講師紹介}

\begin{itemize}
 \item{岩松 信洋} 普段はSuperHですが、今日はポリシーの人です。
 \item{上川 純一} 宴会の幹事です。
\end{itemize}

\subsection{事前課題紹介}

今回の事前課題は
「2006年のDebian勉強会でやりたいこと」
というタイトルで200-800文字程度の文章を書いてください。
というものでした。
その課題に対して下記の内容を提出いただきました。

\subsubsection{岩松 信洋}

\begin{itemize}
 
 \item  ユーザー向けになんかやる\\
  去年は開発者をターゲットにした勉強会だったが、今年は一般ユーザーを対象にした
  ことをやりたい。例えば、Debianでムービーを見るにはどうしたらいいのかなど。
  その前にインストールは全然大変じゃないとこを教えてあげないといけないかも。
  
 \item  出張勉強会をやる\\
  去年は関西オープンソースだけだったけど、OSCとかでもやれたらいいな。

 \item  Debianポリシー\\
  今年から始めるDebianポリシー入門をコンスタントに続けれるようにしたい。


\end{itemize}



\subsubsection{えとーさん}

デベロッパー向けがメインだったDebian勉強会で
ユーザ寄りのことも行なうようになったので、
その中で出来ることをすこしずつやって行きたいと思います。

ユーザ寄りと言っても濃いユーザ寄りの情報交換(or情報発信)の場として
育って行くとうれしいです。
新人発掘のためになにか出来るとよいなと思いますが、まず自分が
新人になるところからでしょうか、デベロッパーはちょっと辛いのですが、
qwikのメンテを通してなんかできるようになりたいなと思います。

\subsubsection{澤田さん}

やりたいこと、debian policyですね。一カ月ぐらい前にdh-make-perlを知って、そ
ういえばPerlのモジュールパッケージって
\verb!lib$(echo $module | tr A-Z a-z | sed -e 's/::/-/g')-perl!
ていう命名規則になってるな〜と思いそこらへんをもっと知りた
くなりました。パッケージ名以外にも、このファイルはどこに置くべき逆にここに置
いてはいけないとか、RecommendsにするかSuggestsにするかの判断基準とか、-commonに分ける判断基準とか。

聞いてばっかりなので発表もしてみたいですね。持ちネタがあまりありませんが。

\subsubsection{吉田＠板橋さん}

2006年のDebian勉強会でやりたいこと
「やってほしいこと」ではなく、「やりたいこと」なところに個人的に
危険な香りを感じますが、それは置いておいて..
.
やってみたいこと、「野良Debianパッケージの作り方」
「使いたいアプリ/機能が公式パッケージにない、複数マシンにインストールし
たり、環境再構築時に楽をしたい、新規パッケージ
はstatbleに入らないからsidとかupstreamから持ってきたいぞ」
といった欲望のままにパッケージを作りたい人のための勉強会
オフィシャルデベロッパーとか英語でITP/RFPが出来る人にはまるで無意味な
勉強会をやってみたい。

\verb!#!オフィシャルデベロッパー向けの勉強会とセットでやるともっと面白いかも...

\subsubsection{中島さん}

　歩行やキーボードを叩くのが困難な身体に障害をもつユーザの方に実際に使用してい
てユーザビリティはどうか。などを教えていただく。テキスト読み上げるブラウザや点
字出力など、そういった障害者用のソフトを集めて使い勝手の良さなどを教えてもらい
たい。　あとどうやってインストールしたのかなど聞かせていただきたい。どのような
使い方、使用目的は何か。などあれば教えていただきたい。キーボードやタッチパネル
はどこのメーカのが良く。どんな入力の仕方をしているのか。そういうことをぜひ聞い
てみたい。



\subsubsection{関根さん}

「隣の芝生は青く見える」

普段から Fedora を使用し、いくつかのゲームの RPM パッケージを作成して
いる立場から見て、Debian がどう見えるのか、Debian の何がよくて何が駄目
かを探っていきたいと思います。逆に、Fedora の素晴らしいところを紹介し、
Debian に活かせるもの伝えていければとも思います。

\subsubsection{上川}

昨年はてさぐりでDebian勉強会を開始しました。昨年の段階では継続して勉強会
をすることがはたしてできるものなのか、というのも明確ではなかったと思いま
す。今年は、勉強会をすること自体は可能だということをふまえて、参加してい
て意味のある内容として継続できるものを模索していきたいと思っています。



%%% trivia quiz
\dancersection{Debian Weekly News trivia quiz}{上川 純一}

ところで、Debian Weekly News (DWN)は読んでいますか？
Debian 界隈でおきていることについて書いているDebian Weekly News.
毎回読んでいるといろいろと分かって来ますが、一人で読んでいても、解説が少
ないので、
意味がわからないところもあるかも知れません。みんなでDWNを読んでみましょう。

漫然と読むだけではおもしろくないので、DWNの記事から出題した以下の質問にこたえてみてください。
後で内容は解説します。

\subsection{2005年50号}
\url{http://www.debian.org/News/weekly/2005/50/}
にある2005年12月13日版です。

\santaku
{www.skolelinux.orgについて提案されていないことはどれか}
{バグトラッキングシステムを共有する}
{関係者のblogをplanetでアグリゲートしたりする}
{こっそりと人を誘拐してメンバーを増やす}
{C}

\santaku
{Branden RobinsonがTuxjournalでのインタビューで、Debianの成功に貢献した
ものとしてあげたのは}
{自由なライセンスを強調する人達とソフトウェアの品質を強調する人達がそれ
ぞれ貢献できてきたこと}
{Rubyをがんがんつかってコードを書いた事}% 彼はpythonユーザです。
{お互いに仲のわるい開発者たちが足をひっぱりあいながらお互いを潰しあって
いたこと}
{A}

\santaku
{GPLでリリースされているゲームボーイ用のエミュレータはmainにいれてもよい
のか}
{フリーのゲームを開発しているグループがあるため、mainにいれてもよい}
{ゲームは商用のゲームしか存在しないためcontribにいれる必要がある}
{エミュレーションという不純な動作はnon-freeにあるべきだ}
{A}

\santaku
{パッケージを実行用のパッケージ {\it pkg} と データ用のパッケージ{\it pkg}-data に分割する場合の確認項目についてBill Allombertは投稿し
た。そこで説明していなかったのは}
{パッケージ名は pkg と pkg-data にしてほしい}
{pkg-data は architecture: all にしてほしい}
{pkg-dataのサイズは5MBを越えていることが望ましい}
{C}

\santaku
{tetexの設定ファイルについてFrankの提案したのは？}
{/usr/share/texmfにデフォルトがあり、/etc/texmf にアドミニストレータの設
定があり、HOME/texmf に各ユーザの設定がある構成}
{Debianメンテナが一番偉いのでユーザの設定を無視して、世界統一の設定にすること}
{Debian menu システムの設定システムが優秀なので、それをそのまま採用する
こと}
{A}



\subsection{2005年51号}
\url{http://www.debian.org/News/weekly/2005/50/}
にある2005年12月20日版です。

\santaku
{debianforum.deは開始何年たったか}
{3年}
{4年}
{5年}
{B}

\santaku
{Jaldhar H. Vyas はインドでは通信コストが高いため、雑誌に付録DVDとしてDebianを
付けたいと提案しました。ただ、複数枚はコストがかかるので、DVD１ 枚におさめたいと
説明しました。Joerg Jaspertの回答は}
{Cebitなどの展示会で利用するためにすでに作成したことがあるので結構簡単だよ、と回答した}
{そんなものつくることがおこがましい}
{ソースだけだったら１枚でもいけるかも}
{A}

\santaku
{lsb向けの起動スクリプトの利用方法について検討していた際に、エラーが発生した
場合にコマンド自体のエラーが画面に出力されて表示がみだれた。
この問題に対して提案された解決策は}
{エラーは/dev/nullへ}
{エラーなんておきないようにする}
{エラーなどをsyslogに送信してみる}
{C}

\santaku
{dpkg-sig を利用してDebianパッケージに電子署名を追加することができる。
最近dpkg-sig の署名を含むDebianパッケージがDebian archiveにアップロードできなくなってい
た。その理由は}
{dpkgにそんな機能拡張はしてはいけないという主義主張の問題}
{予期しない原因でjennifferのチェックが厳しすぎたため}
{実はdpkg-sigなんてものはなかった}
{B}

\santaku
{TexLiveパッケージのライセンスでJoerg Jaspertがおかしいと指摘したのは}
{Liveという名前がダメだ}
{texは時代遅れです}
{DFSGという謎のライセンスを利用している部分が存在した}
{C}

\subsection{2005年52号}
\url{http://www.debian.org/News/weekly/2005/52/}
にある12月27日版です。

\santaku
{Norbert Tretkovskiは、backports.orgで何がおきたと発表したのか。}
{backports.orgがetchに対応した}
{backports.orgのメンテナンスをあきらめた}
{backports.orgがsargeに対応した}
{C}

\santaku
{http://wiki.debian.org/StatusOfUnstable は何を説明してくれるページか}%\url{http://wiki.debian.org/StatusOfUnstable}
{現状のunstableで何がおきているのかをまとめている wikiページ}
{unstableであるということはどういうのかといういことを熱く語るスレッド}
{今どういうことがunstableになりえるのかということを解説しているページ}
{A}

\santaku
{Kevin Lockeが発表したpowermgmtプロジェクトは何をするものか}
{Debianの中での共通の電源管理用のインフラを提供することを目標とする}
{ハックに必要な栄養の補給方法について検討することを目標とする}
{権力をいかに分配するのかということについて検討することを目標とする}
{A}


\subsection{2006年1号}
\url{http://www.debian.org/News/weekly/2006/01/}
にある1月3日版です。

\santaku
{Debian パッケージを圧縮しなおすことで一番小さくできたのはどの圧縮ソフト
ウェアか}
{gzip}
{bzip2}
{7-zip}
{C}

\santaku
{apt-torrentは何をするものか}
{apt風のインタフェースでbittorrentを利用できるツール}
{aptのパッケージダウンロードをbit-torrent経由で実行するためのツール}
{海流予測用ツール}
{B}

\santaku
{vim-tinyは何をするものか}
{nviのかわりにデフォルトにするためにつくられた}
{ただvimを小さくしてみました}
{vimの機能はむだなものが多いので、普通いらないだろうというものだけにして
みた}
{A}


\santaku
{Lars Wirzeniusの提案したのは何か}
{Debianの品質改善のための提案}
{Debianのパッケージ削減のための提案}
{Debianの利用方法の改善のための提案}
{A}

\subsection{2006年2号}
\url{http://www.debian.org/News/weekly/2006/02/}
にある1月10日版です。

\santaku
{Technical Committeeに参加した新メンバーは誰か}
{Steve Langasek, Anthony Towns, と Andreas Barth }
{Wichert Akkerman, Jason Gunthorpe, と Guy Maor }
{Branden Robinson, Kenshi Muto, と Goto Masanori}
{A}


\santaku
{カーネルに存在していたnon-free firmware blobについては現状どうなってい
るか}
{進展がない}
{ライセンスを変更することで全て対処した}
{request firmware というフレームワークによりユーザ空間に移動した}
{C}

\santaku
{apt-get updateでgpgエラーが発生した、これは何か}
{Debianのアーカイブに侵入されたため}
{Debian アーカイブ署名キーが毎年かわるため、2006年用のものに変更する必要があっ
た}
{gpgはもうサポートされていない}
{B}

\subsection{2006年3号}
\url{http://www.debian.org/News/weekly/2006/03/}
にある1月18日版です。

\santaku
{coldfireとはどういうCPUか}
{ヒートシンクに液体冷却を採用したCPUの総称}
{実はCPUではない}
{組み込み用のCPUで、最近はMMUのあるバージョンもある、m68kと一部バイナリ
互換}
{C}

\santaku
{Anthony Townsが宣言したアーカイブの分割とはどういうものか}
{Debian Projectは今後i386のみを配布することに決定した。}
{メインアーカイブはi386, x86\_64, ppcのみにして、その他はscc.debian.org
で提供する}
{メインアーカイブはi386のみにして、それ以外は別のミラーサーバ名で提供す
るようなインフラを準備する}
{C}

\santaku
{協調メンテナンスをするためのフレームワークとして重要だとRaphael Hertzog
が提案したのは何か}
{SVNで管理している協調管理のソースの進捗状況をトラッキングするためのツール}
{仲よくするための定期的な宴会}
{説明書をしっかり書く事}
{A}

\dancersection{最近のDebian関連のミーティング報告}{上川 純一}

\subsection{東京エリアDebian勉強会11回目報告}

% (query-replace-regexp "<.*>" "")
	  12月の第11回Debian勉強会を実施しました。
	    今回は一年間のDebian勉強会をやってみて、の反省会を主としてやりました。
	  今回の参加人数は9人でした。
	  議論された点を以下に紹介します。
        
	  今年のDebian勉強会をふりかえってみました。
	    毎回の参加人数とか、どういうトラブルがあったのか、を振り返りました。
	    また、どういう企画があったのか、ということも説明しました。
	  
	  
	    Debian勉強会のインフラについて説明しました。
	    特に、TeXでのソースがどう作成されているのか、ということを説明しました。
	    hyperrefパッケージでdvipdfmx経由で処理した場合に、URLなどが
	    ハイパーリンクとしてPDFで見れるということに新鮮なおどろきを感じていただけたようでありがたいです。
	    また、クイズ専用マクロとかが回をかさねてどんどん成長しているため、
	    TeX以外で原稿を書くのがどんどん難しくなる、ということを説明しました。
	  
	  
	    今後の企画としては、さらにDebian関連の知識をコアなメンバーの人達にはきだしてもらいたい、ということと、
	    グループワークをたまにはやりたいね、という話がでました。
	    また、現在一人1時間以上の時間をかけて説明しているのですが、
	    それではなく、ライトニングトークのような感じでのプレゼンテーションで軽くできるとよいな、という意見が出ました。
	    他のディストリビューションの人を呼んで、このディストリビューションはここがよいぜ、というのを
	    主張してもらう、という企画もできたらいいな、という意見も出ました。
	    また、今後の企画として、web pageはどうですか、という話をやまねさんに、
	    bootするディストリビューションの開発についてみつかさんに、
	    HCLの話をむとうさんに、
	    Debian JPどうなっているのよという話をgniibeさんにしていただけたらいいなぁ、という話がでました。
	  
	  
	    宴会は「安安」という店で焼肉を食べました。
	    日本語がなかなか通じなくてこまりました。
	    その後、Denniesでデビルズチョコレートサンデー。


\dancersection{Debian Policy 入門 第1回}{岩松 信洋}

\subsection{Debianポリシーとは}
    Debian GNU/Linuxのポリシーです。
    Debian GNU/Linuxとして守るべき方針についてまとめられたものです。
    Debianパッケージの内部構成やオペレーティングシステムとして必要
    な設計部について示されており、ドキュメント化されています。

    これらのドキュメントには debian-policy マニュアルと他の部分に
    ついて補足するサブポリシーマニュアル があります。
    現在のdebian-policy マニュアルバージョンは 3.6.2.2 です。
    毎日議論され、修正が加えられています。


\subsection{debian-policyマニュアルの構成}
    deian-policyマニュアルの構成はどうなっているのか。

    主にDebianパッケージの内容になっています。
    以前はdebian-policyマニュアルとDebian パッケージングマニュアルに
    分かれていたのですが、統合されました (3.2.1.1で統合) 。

    以下にdebian-policyの内容をリストにしてみました。

    \subsubsection{Debianアーカイブ}
\begin{itemize}
        \item Debian Free Software Guidelines (DFSG)とはなにか
        \item main / contrib / non-free セクションの説明および
          各セクションに収録されるパッケージの条件について
        \item Copyrightの問題について
        \item サブセクションについて
        \item パッケージに対するプライオリティについて
\end{itemize}
    \subsubsection{バイナリパッケージについて}

\begin{itemize}
        \item パッケージ名について
        \item パッケージのバージョンについて
            日付に基づいたバージョン番号の付け方

        \item メンテナーのパッケージについて
        \item パッケージの説明について
            パッケージの簡単な説明について
            パッケージの詳細な説明について

        \item パッケージの依存について
        \item バーチャルパッケージ
        \item ベースシステムについて
        \item エッセンシャルなパッケージについて
        \item メンテナースクリプト
\end{itemize}

    \subsubsection{ソースパッケージについて}
\begin{itemize}
        \item 規格への対応
        \item パッケージ関係
        \item 上流パッケージソースの変更について
        \item Debian changelog（debian/changelog）
            代替のchangelog形式
        \item Makefile内でのエラーのトラップについて
        \item タイムスタンプスタンプ
        \item ソースパッケージの中の物における制限
        \item メインビルドスクリプト: (debian/rules)
        \item Variable substitutions: (debian/substvars)
        \item 生成されたパッケージリスト: (debian/file)
\end{itemize}

    \subsubsection{コントロールファイルについて}
\begin{itemize}
        \item コントロールファイルの構文について
        \item ソースパッケージ制御ファイル(debian/control)
        \item バイナリパッケージ制御ファイル(debian/control)
        \item Debianソース制御ファイル--.dsc
        \item Debian Change ファイル--.changes
        \item コントロールファイルのフィールドリスト
        \item ユーザによって定義されたフィールド
\end{itemize}

    \subsubsection{パッケージメンテナンススクリプトとパッケージがインストールされる手順について}
\begin{itemize}
        \item パッケージメンテナスクリプトの序論
        \item メンテナスクリプトの再入結果の同一性
        \item メンテナンススクリプトからのターミナルの制御
        \item メンテナンススクリプトの呼ばれ方の詳細
        \item インストールおよびアップグレードのアンパックフェーズの詳細
        \item 詳細な構成
        \item パッケージの削除とパッケージ設定の完全削除の詳細
\end{itemize}

    \subsubsection{パッケージ同士の関係について}
\begin{itemize}
        \item パッケージ関係フィールドの構文
        \item バイナリの依存について
            （Depends, Recommends, Suggests, Enhances, Pre-Dependsの説明）
        \item バイナリパッケージのコンフリクト( Conflicts )
        \item バーチャルパッケージ( Provides)
        \item ファイルを上書きし、パッケージを置き換える( Replaces )
            他のパッケージの中のファイルを上書きする
            パッケージの削除を強制して、全体のパッケージを置き換える
\end{itemize}

    \subsubsection{共有ライブラリについて}
\begin{itemize}
        \item ldconfig
        \item ランタイムサポートプログラム
        \item スタティックライブラリ
        \item 開発ファイル
        \item 同じライブラリのパッケージとのの依存関係
        \item ライブラリと他のパッケージとの依存 ( shlibsシステム )
            システム上の現在のshlibsファイル
            dpkg-shlibdepsとshlibsファイルの使い方について
            shlibs File フォーマット
            shlibsファイルの提供する
            debian/shlibs.local file を書く
\end{itemize}
    \subsubsection{オペレーティングシステムについて}
\begin{itemize}
        \item ファイルシステム階層構造 ( FHS )
        \item ユーザとグループ
        \item システムランレベルとinit.dスクリプト
        \item init.dスクリプトからのコンソールメッセージ
        \item Cron ジョブ
        \item メニュー
        \item Multimedia handler( MIME )
        \item キーボード構成
        \item 環境変数
        \item doc-baseパッケージを使ったドキュメントの登録方法
\end{itemize}
    \subsubsection{各種ファイルについて}
\begin{itemize}
        \item バイナリファイル
        \item ライブラリファイル
        \item 共有ライブラリ
        \item スクリプト
        \item シンボリックリンク
        \item デバイスファイル
        \item 設定ファイル
        \item ログファイル
        \item パーミッションと所有者
\end{itemize}

    \subsubsection{アプリケーションの変更について}
\begin{itemize}
        \item アーキテクチャ指定のための文字列
        \item デーモン
        \item 仮想tty の使用、wtmp,utmp,lastlog等の更新について
        \item エディタとページャについて
        \item Webサーバーとアプリケーション
        \item メール配送、配信、ユーザーエージェント
        \item ニュースシステムの設定
        \item X Window System 用のプログラム
        \item Emacs Lisp プログラム
        \item ゲーム
\end{itemize}
    \subsubsection{ドキュメントについて}
\begin{itemize}
        \item マニュアル( man pages )
        \item Infoフォーマットのドキュメント
        \item 追加ドキュメント
        \item ドキュメントの管理
        \item 推奨されるドキュメント形式
        \item 著作権関連情報
        \item 設定例
        \item Changelog ファイル
\end{itemize}
    \subsubsection{付録}
\begin{itemize}
        \item Debian パッケージ パッケージングマニュアル
\end{itemize}

\subsection{サブポリシーマニュアルについて}
    メインのものはdebian-policyとして存在し、そのほかにEmacsやPerlに
    関してのサブポリシーマニュアルというものが存在します。
    以下にサブポリシーマニュアルの簡単な説明を書きます。

    \subsubsection{build-essential パッケージの一覧}
        debianのシステム起動に必要なパッケージをリストにしています。
        このパッケージが規定されているドキュメントは

        /usr/share/build-essential/list
 
        になります。
        buiid-essential のリストは

        /usr/share/doc/build-essential/essential-packages-list

        に書かれており、build-essential としてパッケージに収録されています。
        (アーキテクチャによって内容が異なります。）

    \subsubsection{メニューシステム}
        menuシステムを使うためのポリシー。
        引数を持たずに起動可能なアプリケーション(GIMPやxChatなど)はメニューを使って起動で
        きるようにするべきであり、どのようなアプリケーションがどのメニュー項目に入れるべ
        きであるか、書かれています。
        debian-policy パッケージに収録されており、

        /usr/share/doc/debian-policy/menu-policy.txt.gz

        にインストールされます。

    \subsubsection{MIME サポート}
        MIME( Multipurpose Internet Mail Extension RFC1521 )をサポートするためのポリシーです。
        MUAやウェブブラウザでMIMEを扱えるようにできる仕組のようです。
        debian-policy パッケージに収録されており、

        /usr/share/doc/debian-policy/mime-policy.txt.gz

        にインストールされます。

    \subsubsection{Emacs ポリシー}
        Emacs に関連するパッケージは、サブポリシードキュメントに従うことが求められています。
        それをまとめたものが Emacsポリシーです。
        emacsen-common パッケージに収録されており、

        /usr/share/doc/emacsen-common/debian-emacs-policy.gz

        にインストールされます。

    \subsubsection{Java ポリシー}
        Java に関連するパッケージのサブポリシー。
        ドキュメントは
        java-common パッケージに収録されており、

        /usr/share/doc/java-common/debian-java-policy/index.html

        にインストールされます。

    \subsubsection{Ruby ポリシー}
        Ruby に関連するパッケージのサブポリシー。
        ruby パッケージに収録されており、
        ドキュメントは

        /usr/share/doc/ruby/ruby-policy.txt.gz

	にあります。

    \subsubsection{Perl ポリシー}
        Perl に関連するパッケージのサブポリシー。
        debian-policy パッケージに収録されており、

        /usr/share/doc/debian-policy/perl-policy.txt.gz

        にあります。

    \subsubsection{Python ポリシー}
        Python に関連するパッケージのサブポリシー。
        python パッケージに収録されており、

        /usr/share/doc/python/python-policy.txt.gz

        にインストールされます。

    \subsubsection{Debconf 仕様書}
        Debconfのための仕様書。現在プロトコル2。
        debian-policy パッケージに収録されており、
        ドキュメントは

        /usr/share/doc/debian-policy/debconf\_specification.txt.gz

        にあります。

    \subsubsection{スペル辞書・ツールポリシー}
        パッケージの中で使う単語やisspellパッケージやmyspellパッケージのためのポリシー。
        2003年ごろからメンテナンスされてません。
	ドキュメントは

	http://dict-common.alioth.debian.org/

	にあります。
\subsection{どのようにしてポリシーが決まるのか}

    \subsubsection{policy-process}
        /usr/share/doc/debian-policy/policy-process.html/

        にdebianポリシーの決め方が書かれています。

    \subsubsection{debian-policy@list.debian.org (ML)があります}
        ポリシーに関する疑問はこのMLに投げるといいでしょう。
        policy-processに則って、ここで議論され、承認されたときにdebian-policyとして反映されます。

    \subsubsection{debian-policyというパッケージがあります}
        間違いや提案はこのパッケージに対してBTSを行います.
        BTSされたメールはdebian-policy@list.debian.orgにforwardされます。

    \subsubsection{debian-policyのメンテナ}
        現在のdebian-policyのメンテナは以下の4人です。

    \begin{itemize}
        \item Julian Gilbey
            devscripts, tetexのメンテナ
        \item Branden Robinson
            X Strike Froce
        \item Josip Rodin
            debbugs, lintianのメンテナ
        \item Manoj Srivastava
            make, selinuxのメンテナ
    \end{itemize}

\subsection{次回}
	次回からは、debian-policyを一つづつチェックして、つっこんだ解説をしていこうと思っています。


\dancersection{コミケの報告}{岩松 信洋}

Debian勉強会で作成した資料を本にし、コミックマーケットで販売をしました。
以下に報告します。

\subsection{イベントについて}
\begin{itemize}
 \item イベント名

	コミックマーケット６９
 \item 開催日時

	2005年12月30日
 \item 場所

	東京ビックサイト
 \item 出展

	行いませんでした
 \item 委託先

	美紗緒ネットワーク さん

	ありがとうございました。
\end{itemize}
\subsection{本の内容}

	2005年7回勉強会〜2005年10回勉強会の勉強会資料

	2005年7回勉強会〜2005年10回勉強会Debian Weekly News trivia quizおよび答え
\subsection{印刷}
\begin{itemize}
 \item 印刷所

		キンコーズ　市ヶ谷店
 \item 印刷数

		50部
 \item 表紙代

		380円
 \item 中閉じ代

		7500円
 \item コピー代

		18500円
 \item 宅急便

		1500円
 \item 小計

		27800円
 \item 合計
	
		29274円
 \item 消費税 

		1394円
\end{itemize}

\subsection{販売結果}
\begin{itemize}
	\item 販売金額 

		800円

	\item 販売部数
 
		49部（1冊はサンプルとして提出。）

	\item 売上げ
 
		39200円( 800円*49部 )
\end{itemize}
\subsection{次回のコミケ}
手伝ってくださった方々、どうもありがとうございました。
次回も本を作成して販売したいと思っています。
申込書を購入していない（たぶん）と思いますので、委託になると思うので御協力先を探さないといけません。
そのへんを考える必要があります。



% discussion のためのリソースを準備
\dancersection{個人提案課題}{}
\hfill{}{\large 名前} \underline{\hspace{6cm}}

下記の空欄を埋めてください:

{\LARGE 
Debian の (\hspace{5cm})\\
に注目し(\hspace{6cm})\\
を実現するために2006年のDebian勉強会では下記のテーマを実施します
}
\\\\
企画案：
{\LARGE
\begin{enumerate}
 \item 2006年の計画立案
 \item  \underline{(\hspace{10cm})}
 \item  \underline{(\hspace{10cm})}
 \item  \underline{(\hspace{10cm})}
 \item  \underline{(\hspace{10cm})}
 \item  \underline{(\hspace{10cm})}
 \item  \underline{(\hspace{10cm})}
 \item  \underline{(\hspace{10cm})}
 \item  \underline{(\hspace{10cm})}
 \item  \underline{(\hspace{10cm})}
 \item  \underline{(\hspace{10cm})}
 \item 一年間の反省
\end{enumerate}
}

\newpage

%\vfill{}
%\hfill{}
%\includegraphics[width=7cm]{image200502/openlogo-nd.eps}
%\hfill{}
%\vfill{}
%\newpage

\dancersection{グループ提案課題}{}
{\large 名前} \underline{\hspace{6cm}} 
\hfill{}{\large 名前} \underline{\hspace{6cm}}

{\large 名前} \underline{\hspace{6cm}} 
\hfill{}{\large 名前} \underline{\hspace{6cm}}

{\large 名前} \underline{\hspace{6cm}} 
\hfill{}{\large 名前} \underline{\hspace{6cm}}\\

{\LARGE 
Debian の (\hspace{5cm})\\
に注目し(\hspace{6cm})\\
を実現するために2006年のDebian勉強会では下記のテーマを実施します
}
\\\\
下記の空欄を埋めてください:
企画案：
{\LARGE
\begin{enumerate}
 \item 2006年の計画立案
 \item  \underline{(\hspace{10cm})}
 \item  \underline{(\hspace{10cm})}
 \item  \underline{(\hspace{10cm})}
 \item  \underline{(\hspace{10cm})}
 \item  \underline{(\hspace{10cm})}
 \item  \underline{(\hspace{10cm})}
 \item  \underline{(\hspace{10cm})}
 \item  \underline{(\hspace{10cm})}
 \item  \underline{(\hspace{10cm})}
 \item  \underline{(\hspace{10cm})}
 \item 一年間の反省
\end{enumerate}
}

\dancersection{次回}{}

未定です。
内容は本日決定予定です。

参加者募集はまた後程。

\newpage

\vspace*{15cm}
\hrule
\vspace{2mm}
\includegraphics[width=2cm]{image200502/openlogo-nd.eps}
\noindent \Large \bf Debian 勉強会資料\\ \\
\noindent \normalfont 2006年1月21日 \hspace{5mm}  初版第1刷発行\\
\noindent \normalfont 東京エリア Debian 勉強会 （編集・印刷・発行）\\
\hrule

\end{document}
